\label{sec:integer12-saturated}

The polynomial in \path{integer12_saturated/rational.json} has twelve
unrestricted integer unknowns,
123 binary arithmetic operations, and exact total degree 1084.  It
strictly improves both operation count and degree over the pinned
scalar-carrier parents, which have metrics $145/1200$ and $141/1456$
\cite{scalar12}.
The construction combines coprime Pell saturation with the scalar
intersective carrier.  Its completeness requires a rescaling of the
auxiliary coordinate and the full arbitrary-residue interpolation theorem.
The new rescaling and its compatibility with that retained theorem
are proved below.

\subsection{Complete twelve-integer polynomial}

The independent coordinates are
\[
 (j,C,z,h,r,v,w,T,c,n,t,s)\in\Z^{12}.
\]
Use the exact source34 \texttt{mixed-parity-unrounded-1064921} coefficient
tuple, separately from the compact44 tuple of the positive-domain forms:
\[
 \beta=2^{1064921}-1,\qquad
 \kappa_1=2^{1113},\qquad \kappa_2=2^{666104}.
\]
The remaining fixed integers $K_0,M_C,M_F$ are fully specified by the
included six payloads and the finite source34 recipe, complete tile table,
and exact machine seed.  The supplied $M_F$ is already shifted.  Every
coefficient is reconstructed and compared byte for byte by the checker.
No free coefficient parameter remains.

Define the following ordinary polynomial abbreviations; none is another
unknown.  In this section $E$ denotes the first Pell radicand.
\begin{align*}
 J&=6j+1,&q&=\beta J+1,&W&=12w+4,\\
 X&=qW,&Y&=(2v+1)q^3,&a&=(X+1)Y,\\
 H&=4a+3,&\Delta&=a^2+H,&F&=(K_0+W)C-z(q-1).
\end{align*}
\begin{align*}
 R&=\bigl(q(q-F)-C\bigr)(q^2-1)+(M_C+qM_F)J,\\
 k&=R+1+hXY,\qquad \eta=c-Yk,\qquad\zeta=k-\eta,\\
 G&=(\kappa_1F)(\kappa_2C)-(q-\kappa_1F-\kappa_2C)^2,\\
 \Pi&=G\eta\zeta-(t^2-G-\eta\zeta+1)^2,\qquad
 D_g=(3\Pi-4)s^2+1,\\
 \ell&=XY^2,\qquad E=1+\ell(\ell+1)k^2,\qquad U=src^2,\\
 D_f&=1+\Delta U^2E,\qquad Q=\Delta^2U^2E,\qquad\delta=Q-1,\\
 b&=c-R+4UET,\qquad V=D_fb-c,\qquad C_y=QV^2-1,\\
 K&=X(X+2ac)-1,\qquad S=UH^5,\qquad z_c=HT+K,\\
 A_c&=(z_c^2+3H^2)(z_c^3-19H^3),\qquad Z=S(n+V)-A_c.
\end{align*}
Put
\begin{align*}
 f_s&=\delta Z^2,& f_a&=9\delta S^2ED_f,\\
 f_b&=81\delta S^2D_g,& f_c&=S^2C_y,\\
 P_e&=(f_s+f_c-f_a-f_b)^2-4(f_sf_c+f_af_b).
\end{align*}
The entire requested equation over $\Z^{12}$ is
\begin{equation}
 \boxed{\quad P=P_e^2-64f_sf_cf_af_b=0.\quad}
 \label{eq:saturated-integer12}
\end{equation}
Every operation is polynomial arithmetic.  The literal circuit evaluates
$81D_g$ directly as $(243\Pi-324)s^2+81$, saving one multiplication.
Fixed powers in the display are implemented by counted binary
multiplications.

For interpretation, let $\mathcal F_8(Z;A,B,C)$ be the product of
$Z-\epsilon_1\sqrt A-\epsilon_2\sqrt B-\epsilon_3\sqrt C$ over the eight
sign choices.  Sign invariance makes this an integer polynomial.  Then
\begin{equation}
 P=\delta^4\mathcal F_8\left(
 Z;9S^2ED_f,81S^2D_g,\frac{S^2C_y}{\delta}\right).
 \label{eq:saturated-integer12-norm}
\end{equation}
This is only a norm interpretation of the division-free
definition~\eqref{eq:saturated-integer12}.  The denominator is cleared
symbolically, and $\delta$ will be shown never to vanish on integer inputs.

\subsection{Signed bootstrap, separation, and source restoration}

\begin{theorem}[Twelve-integer saturated carrier]
Under the same ordinary native, compiler, and refutation calibration as the
pinned scalar-carrier parent,
\[
 \ZFC\vdash\left(
   \forall x\in\Z^{12}\;P(x)\ne0
 \right)\longleftrightarrow\Con(\ZFC).
\]
The new arithmetic equivalence is proved below and is formalizable in
ZFC.  This statement does not claim a new encoded object-ZFC derivation
or proof-assistant endpoint for the entire inherited calibration.
\end{theorem}

\begin{proof}[Soundness]
The fixed parity and residue choices give, unconditionally,
\[
 |q|\ge16,\quad |X|\ge64,\quad |Y|\ge4096,\quad
 H\ne0,\quad\Delta=(a+2)^2-1>1.
\]
Also $E\ge1$, since $\ell(\ell+1)\ge0$ for integer $\ell$.
If $U=0$, then $Q=0$; otherwise $Q\ge\Delta^2>1$.  Thus $\delta\ne0$
for every integer assignment.

Write $\pi(x)=(x^2+3)(x^3-19)$.  If $U=0$, then $S=0$ and
the final polynomial becomes $A_c^8$.  But
$A_c=H^5\pi(T+K/H)\ne0$: the quadratic has no real zero and the cubic
has no rational zero.  Hence a zero of $P$ has $U\ne0$, so $s,r,c$
are nonzero and $|U|\ge c^2$.  We obtain
$D_f>|c|+1$.  Since $b$ is integral, either $b=0$ and $|V|=|c|\ge1$,
or $|V|\ge D_f-|c|>1$.  Consequently $C_y>0$ and $Q>1$ before any
guard or projection is assumed.

Divide a vanishing conjugate factor of~\eqref{eq:saturated-integer12-norm}
by $S$.  If $D_g<0$, its summand is the only nonzero imaginary term,
an impossibility.  Thus $D_g\ge0$.  Since $s\ne0$, the inequality
$\Pi\le0$ would force $D_g\le-3$; hence $\Pi>0$.  This implies both
$G>0$ and $\eta\zeta>0$.  Indeed their product is positive; if they
were both negative, their absolute values $A,B>0$ would give
$\Pi=AB-(t^2+A+B+1)^2<0$.

Modulo 3, $X=-1$ and $a=H=\Delta=0$.  The value $\ell=XY^2$ is
0 or $-1$, so $E=D_f=1$ and $\delta=C_y=-1$ modulo 3.  The three
normalized roots have squared 3-adic orders
\[
 \sqrt{C_y/\delta}:0,\qquad
 3\sqrt{ED_f}:2,\qquad
 9\sqrt{D_g}:4+v_3(D_g).
\]
Omit a zero guard term.  In a rational sum of square roots of positive
rationals, each nontrivial square class must cancel internally.  Within
such a class the ratios are rational, and the displayed valuations
give a unique term of minimum 3-adic order, making cancellation
impossible.  Therefore all three roots are rational.

The integer $ED_f$ is thus a square.  Since
$D_f=1+\Delta U^2E$, we have $\gcd(E,D_f)=1$; positivity and unique
factorization give
\[
 E=\tau^2,\qquad D_f=f^2,\qquad I=U\tau,
 \qquad \tau,f\in\Z_{>0}.
\]
This includes $E=1$.  We have $c^2\mid I$, $I\ne0$,
$f^2=1+\Delta I^2$, and $Q=(\Delta I)^2$.
The rational auxiliary root $y$ satisfies
\begin{equation}
 QV^2-(Q-1)y^2=1.
 \label{eq:sat12-aux}
\end{equation}
The adjacent-coefficient descent from the retained parent makes $y$
integral.  Explicitly, if $|y|=a/b$ in lowest terms then $b^2\mid Q-1$;
the norm-preserving map
\[
 V'=(2Q-1)V-2(Q-1)y,\qquad y'=(2Q-1)y-2QV
\]
preserves the reduced denominator and strictly decreases $|V|$ until
$|V|=1$, where $|y|=1$.  Signs are reflected during the descent.
The guard root is integral as well.

The vanishing factor now gives
\[
 S(n+V)-A_c=S(\pm3\tau f\pm9g\pm y),
\]
so $S\mid A_c$ and $H^5\mid A_c$.  Expanding
\[
 \Phi_H(z)=z^5+3H^2z^3-19H^3z^2-57H^5,
\]
if $p^e\Vert H$ and $v_p(z)=b<e$, its leading term has valuation
$5b$ and every other term has larger valuation.  Thus
$H^5\mid\Phi_H(z)$ implies $H\mid z$.  Applying this to $z=HT+K$
recovers $H\mid K$, with no sign assumption on $H$.

The exact source coordinates are
\begin{equation}
 r_{\rm old}=r\tau,\qquad T_{\rm old}=\tau T,
 \qquad I=sr_{\rm old}c^2.
 \label{eq:sat12-restoration}
\end{equation}
They give
\[
 4IT_{\rm old}=4UE T,\qquad
 V=f^2(c-R+4IT_{\rm old})-c.
\]
All strict guards, the first and strong norm squares,
the auxiliary norm~\eqref{eq:sat12-aux}, $c^2\mid I$, and $H\mid K$
are now exactly the hypotheses of the pinned signed scalar-carrier
restoration theorem.  That theorem permits arbitrary nonzero strong
multipliers, including even ones.

For precision, it uses absolute Pell folding and Lucas multiplication
to recover $|c|=\psi_{|a+2|}(R)$; the auxiliary residues and strict
index bounds identify the index.  Its whole-product estimate, using
$H\mid K$, then yields $\operatorname{sign}(c)X=2^R$ and restores
the missing main root and generalized signed native kernel.  These
unchanged lemmas are retained in \path{integer12_saturated/parent/PROOF.md}
and its pinned signed proof sources.  The named native/refutation
calibration therefore supplies a finite syntactic ZFC refutation.
The claim is restoration of that kernel, not a pointwise substitution
in the old scalar polynomial, whose carrier uses $T_{\rm old}$.
\end{proof}

\subsection{Completeness with a fixed strong witness}

Start with the canonical native witness of a finite ZFC refutation.
Its relevant conditions are
\[
 A_0=a+2>1\text{ even},\quad R\equiv3\pmod4,\quad
 c=\psi_{A_0}(R)>0,\quad c\equiv R\pmod4,
 \quad E=\tau^2,\quad\tau>0.
\]
Choose $t$ for the joint strict guard, then a positive $s$ with
$D_g=g^2$, as in the parent.  When $\Pi=1$ one can use $s=1,g=0$;
otherwise ordinary Pell existence supplies the guard witness.

The full-residue interpolation theorem in the retained
\path{integer12_saturated/parent/source/INTERPOLATION12.md}, Section 6, says that for
\emph{every} prescribed positive $N$, there is one fixed strong pair
$I,f$ with $4Nc^2\mid I$, such that \emph{every} residue class modulo
that same $I$ is attained by an auxiliary $T_{\rm old}$ satisfying the
norm.  Its strong pair is chosen before the residue request, and its
squared-modulus period is $4mf$, where $m$ is the strong Pell index.

Apply this theorem with $N=s\tau$, obtaining
\[
 4s\tau c^2\mid I,\qquad U=I/\tau,
 \qquad r=I/(s\tau c^2).
\]
Thus $U=src^2$ is integral.  The canonical main projection gives
$H\mid K$; put $\kappa=K/H$.  The retained intersectivity theorem for
$\pi(x)=(x^2+3)(x^3-19)$ provides a residue $a_*$ modulo $U$ with
\[
 \pi(a_*+\kappa)\equiv0\pmod U.
\]
Now request the old auxiliary residue
\[
 T_{\rm old}\equiv\tau a_*\pmod I.
\]
Since $I=\tau U$, this ensures $\tau\mid T_{\rm old}$ and, on setting
$T=T_{\rm old}/\tau$, gives exactly $T\equiv a_*\pmod U$.
The new auxiliary expression is therefore the required source
expression, and
\[
 \frac{A_c}{S}=\frac{\pi(T+\kappa)}{U}\in\Z.
\]
Finally choose
\[
 n=\frac{A_c}{S}+3\tau f+9g+y-V\in\Z.
\]
This makes a conjugate factor vanish and proves completeness.
No positivity of $n$, coprimality of $\tau$ and $U$, or change of strong
pair after the residue choice is needed.  The factor $E$ in $4UET$ is
essential to this exact quotient-coordinate argument.

\subsection{Attained degree and reproducible cost}

All fixed coefficients have degree zero.  The relevant total degrees are
\[
\begin{array}{c|rrrrrrrrrr}
 &U&E&\Delta&D_f&\delta&b&V&C_y&S&A_c\\\hline
 \deg&4&34&12&54&66&39&93&252&34&45.
\end{array}
\]
Put $z_0=Sn-A_c$, of degree 45, and $\omega=f_s-f_c$.
The defining relation $C_y=(\delta+1)V^2-1$ gives the exact cancellation
\[
 \omega=\delta z_0^2+2\delta SVz_0-S^2(V^2-1).
\]
Its unique top term is $-S_0^2V_0^2$, of degree 254; competing degrees
are 156, 238, and 68.  The four squared arguments have degrees
\[
 \deg f_a=222,\qquad\deg f_b=180,\qquad
 \deg f_c=\deg f_s=320.
\]
The identity
\[
 P_e=\omega^2-2(f_a+f_b)(2f_c+\omega)+(f_a-f_b)^2
\]
has unique top term $-4(f_a)_0(f_c)_0$, of degree 542; every competitor
has degree at most 508.  The correction $64f_sf_cf_af_b$ has degree
1042, so the final degree is $2\cdot542=1084$, provided the top form
does not vanish at the fixed tuple.

It does not.  With subscript zero denoting highest homogeneous part,
\begin{align*}
 q_0&=6\beta j,&X_0&=12q_0w,&Y_0&=2q_0^3v,\\
 a_0&=X_0Y_0,&U_0&=src^2,&E_0&=h^2X_0^4Y_0^6,\\
 S_0&=2^{10}a_0^5U_0,&\delta_0&=a_0^4U_0^2E_0,
   &(D_f)_0&=a_0^2U_0^2E_0,\\
 V_0&=4a_0^2U_0^3E_0^2T.&&&&
\end{align*}
The term $4UET$ uniquely dominates $R$ in $b$, so these are single
monomials.  The entire degree-1084 homogeneous part is
\begin{equation}
 P_0=2^{1144}3^{760}\beta^{624}
 j^{624}w^{132}v^{164}h^{32}s^{32}r^{32}c^{64}T^4\ne0.
 \label{eq:sat12-top}
\end{equation}
The exponents sum to 1084 and the actual fixed $\beta$ is positive.
Thus the degree is proved for the complete polynomial, independently
of finite line tests.

The literal operation count is 75 multiplications, 25 additions, and
23 subtractions: 123 in total, including scalar multiplications.  Every
arithmetic node is live.  The 13-operation final norm follows a
110-operation prefix and carrier.  An included integral control has the
same operation count and degree 1280; it is dominated by this candidate.

The independent standard-library checker imports no producer.  It
verifies every prefix by free-ring identities, all 35 coefficients of
the eight-conjugate norm, the centered cancellation and restoration
identities, exact leading induction, complete coefficient bytes and
recipe, and operation counts.  Normal and optimized runs pass; corrupted
auxiliary, strong, guard, carrier, and final expressions are rejected.
These computations certify the literal polynomial and its algebraic
metrics.  The signed semantic argument and interpolation proof are
ordinary mathematical proofs using the explicitly retained source
interfaces, independently reviewed in the accompanying audit.
